\section{Appendix: Critical Values and Inferencefor the Cointegration Diagnostics}
\label{app:critical_values}

Three families of test statistic carry the cointegration verdict in this
chapter, and they do not draw their critical values from the same source.
The residual-based diagnostics use published constants. The fixed-$b$ Wald
statistics do not, because their limiting distribution is jointly indexed
by the bandwidth ratio $b$ and the number of integrated regressors $k$, so
no single tabulated value covers the specifications estimated here.

\begin{table}[H]
	\centering
	\caption{Critical Values for the Residual-Based Cointegration Diagnostics}
	\label{tab:critical_values}
	\small
	\begin{tabular}{p{5.6cm}p{4.4cm}p{1.4cm}p{1.4cm}p{1.4cm}}
		\toprule
		Test & Source & 10\% & 5\% & 1\% \\
		\midrule
		Residual ADF ($t_{\text{ADF}}$) & \citet{MacKinnon2010}, Case 2,
		constant only, $m=4$ & $-3.81$ & $-4.10$ & $-4.64$ \\
		Shin LM ($C_{\text{Shin}}$), $k=2$ & \citet{Shin1994}, Table 1,
		demeaned $C_\mu$ & $0.163$ & $0.221$ & $0.380$ \\
		Shin LM ($C_{\text{Shin}}$), $k=3$ & \citet{Shin1994}, Table 1,
		demeaned $C_\mu$ & $0.121$ & $0.159$ & $0.271$ \\
		\bottomrule
		\multicolumn{5}{p{14.2cm}}{\footnotesize \textit{Notes:} The
			residual ADF statistic tests the null of no cointegration, so
			rejection supports a cointegrating relation. The Shin LM
			statistic reverses the null: it tests the null of
			cointegration, so a value below the critical value fails to
			reject and is the favorable outcome. $m$ denotes the number
			of regressors entering the MacKinnon response surface; $k$
			denotes the number of integrated regressors entering the Shin
			statistic.} \\
	\end{tabular}
\end{table}

For the fixed-$b$ Wald statistics reported as $CT_{\text{IM}}$ and for the
coefficient significance stars in every results table, critical values are
generated by simulation rather than read from a table. The protocol holds
the estimated specification fixed and varies only the innovation draws.
For each specification I generate integrated regressors as independent
random walks of the same length as the estimation sample, impose the null
under test, apply the same Frisch-Waugh-Lovell orthogonalization and the
same deterministic components used in estimation, and compute the IM-OLS
statistic on each replication. The empirical quantiles of the resulting
distribution supply the thresholds. Replication counts are 100{,}000 for
the United States specifications and 10{,}000 for the Chilean threshold
models, where each pass carries the additional cost of a profile search
over the threshold parameter.

Two features of this protocol matter for reading the tables. First,
because the simulated null reproduces the bandwidth ratio actually used
in estimation, the reported stars already incorporate the fixed-$b$
correction and are not comparable to standard normal or chi-squared
thresholds. Second, because the Chilean critical bounds are simulated
under profile search, they are conservative relative to a fixed-threshold
benchmark: the search itself consumes degrees of freedom, and the
simulation charges for it.
